\chapter{Feature categoriche ed embedding}
\label{cap:embedding}
\fonti{\rif{2026-10-19}{19/10}, \rif{2026-10-20}{20/10}}{4-embedding-introduction, pp.\ 1--51}

Tutte le architetture viste finora (RNN nel cap.~\ref{cap:rnn}, GNN nel cap.~\ref{cap:cnn-gnn}) servono a \textbf{estrarre una
rappresentazione} dei dati quando costruirla a mano è difficile (le feature di una serie o di un grafo). Questo capitolo affronta lo
stesso problema per un tipo di dato onnipresente in cybersecurity: le \textbf{feature categoriche}. Prima si vedono le codifiche
classiche (ID sequenziali, feature engineering, one-hot), poi gli \textbf{embedding appresi} e i due modi di impararli: con le label
(supervisionati) o senza (self-supervised).

% =====================================================================
\section{Feature categoriche: dalle categorie ai numeri}
Le reti (e molti algoritmi ML) richiedono input numerici e assumono relazioni matematiche significative tra i numeri.
L'obiettivo è convertire le categorie in numeri \textbf{significativi per il problema} e per la rete (assegnare numeri a caso
come 10, 5, 6, 1 non lo è). Le feature categoriche possono essere stringhe o numeri.
\begin{itemize}
  \item Esempi: titoli di lavoro, processi chiamati durante l'esecuzione di un file, algoritmi di cifratura;
        soddisfazione del cliente, livello di accesso a un server, gravità di un attacco.
  \item Per le \textbf{label} la soluzione migliore è quasi sempre l'ID sequenziale (label encoding).
\end{itemize}

\subsection{Prima domanda: le categorie hanno un ordine intrinseco?}
\textbf{Ordinali} (soddisfazione, livello di accesso, gravità): \textbf{ordina e mappa}.
\begin{enumerate}
  \item identifica i valori unici;
  \item ordinali con la \textbf{conoscenza del dominio} (non in ordine alfabetico!) e assegna indici: Low = 1, Medium = 2, High = 3;
  \item mappa il dataset originale.
\end{enumerate}
Il modello non sa cosa significhi la categoria né se ``alto'' sia buono (gravità: basso è meglio; soddisfazione: alto è meglio):
sa solo che 1 è più vicino a 2 che a 3, e lo scoprirà durante il training.

% =====================================================================
\section{Feature nominali: le alternative}
Esempio: AES, DES, RC4, RSA. Non c'è una mappatura intuitiva. Il corso presenta quattro opzioni.

\subsection{1. ID sequenziale comunque}
AES = 1, DES = 2, RC4 = 3, RSA = 4.
\begin{itemize}
  \item Mettiti nei panni del modello: perché AES è più vicino a DES che a RSA? Il ML scopre pattern nei dati: una mappatura sbagliata lo
        mette su una \textbf{falsa pista} (gli dai un \textbf{bias} da ignorare).
  \item Anche se da esperto ordino per sicurezza (meno sicuro = 1, più sicuro = 4): se domani RSA viene rotto, devo cambiare la mappatura
        e il suo significato. Il modello impara ``RSA viene dopo AES'', non perché siano simili.
\end{itemize}

\subsection{2. Feature engineering}
AES, RSA = sicuri = 1; DES, RC4 = insicuri = 0.
\begin{itemize}
  \item \textbf{Pro}: riduce lo sforzo di interpretazione: dai al modello una chiave di lettura, con un vincolo più debole dell'ordine.
  \item \textbf{Contro}: e se la chiave è sbagliata o cambia? (RSA oggi sicuro, domani rotto.) In un problema benigno/malevolo il modello
        potrebbe associare gli algoritmi insicuri al traffico malevolo (bias). E a volte il feature engineering è difficile o impossibile.
\end{itemize}

\subsection{3. One-hot encoding}
\begin{enumerate}
  \item si identifica l'insieme delle categorie, il \textbf{vocabolario} $V$;
  \item ogni categoria diventa un vettore \textbf{sparso} $V$-dimensionale (tutto 0 tranne un 1); i vettori sono \textbf{ortogonali}.
\end{enumerate}
\begin{itemize}
  \item \textbf{Pro}: semplice, \textbf{nessun bias}: tutti i vettori sono equidistanti; con algoritmi tradizionali (random forest) diventa
        $V$ feature binarie (presenza/assenza di AES, di DES\ldots). Funziona anche con lo shallow learning.
  \item \textbf{Contro}: \textbf{maledizione della dimensionalità} (tante feature = problema più difficile; 171.000 parole inglesi);
        se il \textbf{vocabolario si estende} (es.\ ChaCha20, o un futuro algoritmo post-quantum) cambia la rappresentazione e il training fatto
        va buttato: si riparte da zero.
\end{itemize}
\begin{verifica}
Distanza tra vettori one-hot: due vettori one-hot distinti hanno distanza euclidea $\sqrt{2}$ (la distanza \emph{al quadrato} è 2).
La slide scrive ``= 2''.
\end{verifica}

\subsection{4. Embedding appresi (solo reti neurali)}
\begin{enumerate}
  \item si identifica il vocabolario $V$ (es.\ $V=4$);
  \item ogni categoria si mappa su un vettore \textbf{denso} $H$-dimensionale, con $H \ll V$; ogni categoria ha il suo vettore
        ($V$ vettori diversi); all'inizio i vettori sono \textbf{casuali}, senza relazioni;
  \item gli embedding si \textbf{imparano risolvendo un task} (es.\ sicuro vs insicuro) tramite backpropagation, per molte epoch.
\end{enumerate}
Si può fare solo con le reti neurali perché solo lì si può mandare indietro il segnale e cambiare i valori; random forest e simili no.

\begin{esempio}[cosa succede ``dietro le quinte'']
Rete: \code{nn.Linear(V, H)} $\to$ \code{nn.Linear(H, 2)} $\to$ softmax, con $V=4$, $H=2$; input = one-hot degli algoritmi; label = sicuro/insicuro.
\begin{itemize}
  \item \textbf{Epoch 1}: pesi casuali. Dopo il primo layer ogni algoritmo ha una rappresentazione 2D senza senso (es.\ AES $=[-0{,}7, 0{,}3]$).
        Il modello dà probabilità alte a predizioni sbagliate (è sicuro di sbagliare): 3 giuste su 5.
  \item \textbf{Epoch 2}: la backprop cambia i pesi del primo layer, quindi la rappresentazione: lo spazio inizia a separarsi; 4 giuste su 5.
  \item \textbf{Epoch 3}: tutte giuste; nello spazio 2D gli algoritmi sicuri hanno coordinate positive (AES $[2{,}1, 2{,}5]$, RSA $[1{,}9, 2{,}0]$),
        quelli insicuri negative (DES $[-2{,}1, -1{,}9]$, RC4 $[-1{,}1, -1{,}3]$).
\end{itemize}
Il vettore $H$-dimensionale ``incorpora'' (\en{embeds}) le caratteristiche più salienti dell'oggetto in numeri significativi per il task:
per questo si chiama \textbf{embedding}. Si possono impilare più layer per rappresentazioni più astratte.
\end{esempio}

\begin{center}
\small
\begin{tabularx}{\textwidth}{X X}
\toprule
\textbf{Benefici degli embedding} & \textbf{Problemi}\\
\midrule
Rappresentazione \textbf{densa} per costruzione; parte casuale e impara una distribuzione significativa (che ottimizza il task) & Solo con reti neurali (non con random forest)\\
Più efficiente del one-hot: l'informazione non è sparsa & Serve un \textbf{task} che dia il segnale di backprop (supervisionato o self-supervised)\\
Le feature codificano \textbf{semantica}, non solo informazione binaria & Se il vocabolario si estende, la lookup ha ancora $V$ righe: ri-addestrare da zero? training incrementale?\\
Dimensione molto ridotta ($H \ll V$: es.\ 128 invece di 171.000) & \\
\bottomrule
\end{tabularx}
\end{center}

\section{Categorie sconosciute: il token \code{<UNK>}}
Un nuovo algoritmo o un nuovo dominio non sono nel vocabolario: non hanno una rappresentazione. Si aggiunge al vocabolario una categoria
\textbf{sconosciuta} (\code{<UNK>}) a cui si mappa tutto ciò che non è stato visto in training. Succede sempre: validation e test possono
contenere valori mai visti (il test set spesso è raccolto molto dopo).

\section{La roadmap: dalle categorie ai numeri}
\begin{center}
\begin{tikzpicture}[node distance=0.55cm and 0.9cm, every node/.style={font=\small, align=center},
  q/.style={draw=polblu, fill=polblu!6, rounded corners=2pt, minimum width=3.4cm, minimum height=0.8cm},
  a/.style={draw=esecol, fill=esecol!10, rounded corners=2pt, minimum width=2.8cm, minimum height=0.8cm}]
  \node[q] (q0) {Feature categorica};
  \node[q, below=of q0] (q1) {Ha un ordine intrinseco?};
  \node[a, right=of q1] (a1) {ID sequenziali};
  \node[q, below=of q1] (q2) {Poche categorie?};
  \node[a, right=of q2] (a2) {One-hot encoding};
  \node[q, below=of q2] (q3) {Posso creare un task\\per una rete neurale?};
  \node[a, right=of q3] (a3) {Embedding};
  \node[a, below=of q3] (a4) {Feature engineering};
  \draw[-{Stealth}] (q0) -- (q1);
  \draw[-{Stealth}] (q1) -- node[above, font=\scriptsize]{sì} (a1);
  \draw[-{Stealth}] (q1) -- node[left, font=\scriptsize]{no} (q2);
  \draw[-{Stealth}] (q2) -- node[above, font=\scriptsize]{sì} (a2);
  \draw[-{Stealth}] (q2) -- node[left, font=\scriptsize]{no} (q3);
  \draw[-{Stealth}] (q3) -- node[above, font=\scriptsize]{sì} (a3);
  \draw[-{Stealth}] (q3) -- node[left, font=\scriptsize]{no} (a4);
\end{tikzpicture}
\end{center}
\begin{itemize}
  \item \textbf{Feature engineering} quando il resto non va: con Intesa Sanpaolo, per 120--1000 transazioni per utente al mese, si sono
        progettate oltre 100 feature riassuntive, in numero gestibile qualunque fosse il numero di transazioni.
  \item \textbf{Si può usare un ID sequenziale per una feature nominale?} Non si \emph{dovrebbe}, ma \textbf{si può}, se sai cosa stai
        facendo e l'ordine ha senso per il tuo task. Non è una scienza esatta: le scelte sub-ottimali si pagano (più dati, più epoch,
        training più instabile) e il rischio è introdurre bias.
  \item La lista non è esaustiva (es.\ \en{feature hashing}).
\end{itemize}

% =====================================================================
\section{Come imparare gli embedding (1): supervisionati}
Gli embedding si possono imparare in due modi (slide 4, p.~40): con un segnale \textbf{supervisionato} o \textbf{self-supervised}.
L'idea: invece di fare noi feature engineering (a volte impossibile), lasciamo che sia il \textbf{task} a guidare la rappresentazione
delle feature categoriche.

\textbf{Ipotesi}: scenario \textbf{supervisionato} (label disponibili: traffico benigno/malevolo, email di phishing, famiglia di botnet\ldots),
dati con feature numeriche e categoriche, e una \textbf{rete neurale}: solo le reti possono imparare embedding. Con un random forest si possono
risolvere classificazione e regressione, ma non fare \textbf{feature extraction}: ``ottenere embedding'' = feature extraction.

\begin{definizione}[lookup table e \code{nn.Embedding}]
\code{torch.nn.Embedding(V, H)} crea una tabella \textbf{addestrabile} con $V$ righe (numero di categorie, la dimensione del vocabolario) e
$H$ colonne (la dimensione dell'embedding, un \textbf{iperparametro}). È equivalente a un one-hot seguito da un layer lineare. Permettendo la
backpropagation, la rete cambia i valori della tabella per ottimizzare il task.
\end{definizione}
\begin{integrazione}
In pratica \code{nn.Embedding} riceve \textbf{indici interi} (\code{torch.long}) e restituisce la riga corrispondente: è il prodotto one-hot
$\times$ matrice dei pesi senza costruire il vettore one-hot, molto più efficiente in memoria.
\end{integrazione}
Si addestra la rete sul task, si aggiustano i pesi dei layer in base agli errori, e la rappresentazione si ``taglia'' a metà rete (dopo il primo
layer e prima dell'ultimo).

\subsection{Esempio 1: dataset tabellare con feature categoriche e numeriche}
Metadati di email (dominio del mittente, numero di link, numero di allegati, lunghezza dell'oggetto, ora di ricezione); task: phishing (1)
o normale (0). Il \textbf{dominio del mittente} ha cardinalità altissima (migliaia di domini con un mail server) $\Rightarrow$ il one-hot creerebbe
milioni di colonne, e non c'è un ordine (Gmail prima di Outlook o di Polito?). $\Rightarrow$ \textbf{embedding}.

\begin{enumerate}
  \setcounter{enumi}{-1}
  \item \textbf{Split} train/validation/test.
  \item \textbf{Vocabolario dal solo training}: è ciò che si suppone di conoscere nella realtà. Il test replica il futuro, e anche la validation
        non va considerata nota. Si aggiunge \code{<UNK>} per i valori mai visti (un dominio nuovo come \code{updatebanking.com} o \code{polito.it} in
        validation $\Rightarrow$ entrambi \code{<UNK>}, stesso embedding). Gli ID servono solo come indici per la lookup, non come feature.
  \item \textbf{Mappare} la feature categorica negli ID.
  \item \textbf{Dataset} con feature categoriche e numeriche \textbf{separate} (vanno trattate diversamente), convertite in tensori.
  \item \textbf{Modello}: la categorica passa nell'\code{nn.Embedding}, l'embedding si \textbf{concatena} alle numeriche, il risultato entra nella FFNN.
        Il primo layer lineare ha \textbf{embed\_dim + num\_numeric} neuroni di input: una sola feature categorica diventa 16, 32\ldots numeri.
\end{enumerate}

\begin{lstlisting}
def build_vocab(train_values):
    vocab = {'<PAD>': 0, '<UNK>': 1}
    for v in train_values:                 # SOLO training set
        if v not in vocab:
            vocab[v] = len(vocab)
    return vocab

def encode(values, vocab):
    return [vocab.get(v, vocab['<UNK>']) for v in values]   # mai visto -> UNK

class EmbFFNN(nn.Module):
    def __init__(self, vocab_size, embed_dim, num_numeric, num_classes=2):
        super().__init__()
        self.emb = nn.Embedding(vocab_size, embed_dim)
        self.net = nn.Sequential(
            nn.Linear(embed_dim + num_numeric, 64), nn.ReLU(),
            nn.Linear(64, num_classes))
    def forward(self, cat_ids, num_feats):
        e = self.emb(cat_ids)                    # (batch, embed_dim)
        x = torch.cat([e, num_feats], dim=1)     # concatena con le numeriche
        return self.net(x)
\end{lstlisting}
L'embedding layer è un layer come gli altri: in backpropagation il segnale arriva anche a lui e ne migliora la rappresentazione.

\subsection{Esempio 2: sequenze di valori categorici}
Una \textbf{darknet} è un insieme di indirizzi IP passivi che non ospitano servizi: riceve per definizione solo \textbf{traffico non richiesto}.
Input: sequenza degli IP che ogni ora contattano dei server (lista di lunghezza variabile); task: il server è stato attaccato da una botnet (e
quale)? Una sequenza $\Rightarrow$ \textbf{RNN} (cap.~\ref{cap:rnn}), in cui ogni elemento è una categoria (un IP). Lo stesso schema servirà
nel Lab 2 con le sequenze di API call (cap.~\ref{cap:lab2}).
\begin{enumerate}
  \item split; vocabolario degli IP dal training; mappatura in ID (come prima);
  \item DataLoader per RNN con padding dentro il batch (o padding dinamico con \code{collate\_fn});
  \item modello: \code{nn.Embedding} sulla sequenza (PyTorch gestisce i singoli elementi), poi LSTM, poi head lineare. L'input della LSTM non è 1
        (una categoria) ma \textbf{embed\_dim}. Si può indicare l'ID del padding (\code{padding\_idx}): il suo embedding \textbf{non riceve backpropagation}.
\end{enumerate}
\begin{lstlisting}
class IPSeqLSTM(nn.Module):
    def __init__(self, vocab_size, embed_dim, hidden, num_classes, pad_idx=0):
        super().__init__()
        self.emb = nn.Embedding(vocab_size, embed_dim, padding_idx=pad_idx)
        self.lstm = nn.LSTM(embed_dim, hidden, batch_first=True)  # input = embed_dim
        self.fc = nn.Linear(hidden, num_classes)
    def forward(self, ids, lengths):
        e = self.emb(ids)                                  # (batch, seq, embed_dim)
        packed = pack_padded_sequence(e, lengths.cpu(), batch_first=True,
                                      enforce_sorted=False)
        _, (h, _) = self.lstm(packed)
        return self.fc(h[-1])
\end{lstlisting}

\subsection{Esempio 3: architetture composte (movimento laterale)}
La feature extraction si può fare con \textbf{qualsiasi} architettura, anche combinandole. Problema: distinguere i \textbf{movimenti laterali}
(spostarsi da una macchina all'altra, es.\ verso il computer di un amministratore) in una rete aziendale.
\begin{itemize}
  \item I \textbf{nodi} sono i servizi attivi nella rete (identificatori categorici).
  \item Ogni \textbf{grafo} codifica i movimenti di un utente sui nodi in un periodo (es.\ 1 ora).
  \item Gli attacchi possono essere \textbf{APT} (\en{Advanced Persistent Threat}): si monitora a snapshot diversi $\Rightarrow$ ogni utente ha una
        \textbf{serie temporale di grafi}.
\end{itemize}
Architettura finale:
\begin{enumerate}
  \item \textbf{embedding layer} (lookup): mappa i servizi in embedding;
  \item \textbf{GNN} (cap.~\ref{cap:cnn-gnn}): rappresenta come l'utente si è mosso tra i servizi in uno snapshot (la stessa GNN per tutti gli snapshot);
  \item \textbf{RNN}: tiene conto dei diversi snapshot e produce la predizione (movimento laterale o attività normale).
\end{enumerate}
Tutto si addestra con \textbf{lo stesso segnale supervisionato}: l'errore della predizione si backpropaga attraverso RNN (nel tempo), GNN (per
nodo) e fino all'embedding layer, che riceve la somma dei gradienti di tutti i rami (timestamp e nodi). Non è magia: è moltiplicazione di matrici.
Un utente legittimo può fare SSH da una macchina all'altra per arrivare al server (passando dalla VPN): non tutto il movimento è malevolo.

\begin{esame}[gli embedding supervisionati sono specifici del task]
\textbf{Si potevano progettare feature a mano?} Molto difficile: servirebbero più feature per ogni IP e servizio (importanza nell'azienda, cosa
fa\ldots), senza sapere se sono adatte. Gli embedding invece sono per costruzione \textbf{adatti a questo problema}.

\textbf{E se domani voglio usare gli stessi servizi per prevedere quale servizio sarà più usato?} Numericamente posso riusare gli embedding, ma
sono stati appresi per un task diverso: \textbf{ottimi per questo task, sub-ottimi per altri}. Con le label, la rappresentazione è funzionale al
problema finale, non generica.
\end{esame}

\subsection{Il rischio: cosa sta imparando? (husky vs lupo)}
Come distingui un husky da un lupo nelle foto? La via più ``pigra'': gli husky sono fotografati sull'\textbf{erba}, i lupi sulla \textbf{neve}.
Il classificatore impara lo \textbf{sfondo}: una \textbf{scorciatoia} (\en{shortcut}) che funziona sui dati di training ma non è generale.
La rappresentazione è ottima per \emph{quelle} label e \emph{quei} dati. È lo stesso fenomeno visto nel Lab 1 e nel discorso sulle feature che
creano bias (porta 80, cap.~\ref{cap:ffnn}).

% =====================================================================
\section{Come imparare gli embedding (2): self-supervised}
\begin{definizione}[self-supervised learning, SSL --- slide 4, p.~43]
Un algoritmo self-supervised \textbf{genera un segnale supervisionato da dati non etichettati}: sono i dati stessi a fornire la supervisione.
È un sottoinsieme dell'apprendimento non supervisionato, ma mentre l'unsupervised \emph{descrive} i dati, il self-supervised usa il segnale
per imparare a \emph{rappresentarli}.
\begin{itemize}
  \item \textbf{Pretext task}: si impara una rappresentazione significativa dai dati non etichettati (es.\ un encoder che comprime l'input
        in uno spazio a dimensione minore, il \en{bottleneck}: è la fase di feature extraction).
  \item \textbf{Downstream task}: si usano quelle feature per il compito vero (classificazione, ricostruzione\ldots).
\end{itemize}
\end{definizione}
Il self-supervised learning è la pietra angolare dei successi recenti del deep learning:
\begin{itemize}
  \item permette di imparare feature \textbf{generiche} (``cosa rende una banana una banana'');
  \item oggi è preferito alle feature end-to-end supervisionate;
  \item \textbf{pretext}: si pre-addestra il modello in self-supervision su grandi corpora;
        \textbf{downstream}: lo si fa fine-tuning con poche label;
  \item cruciale in cybersecurity, dove le label sono poche e si vuole generalizzare.
\end{itemize}

Nel supervisionato il segnale viene dalle label: scoiattolo vs tigre $\to$ feature extractor $\to$ embedding $\to$ classification head
$\to$ predizione sbagliata $\to$ backprop. Senza label: \textbf{da dove viene il segnale?}
Idea generale: sfruttare le \textbf{proprietà interne degli oggetti} per capire cosa è simile e cosa no.

\subsection{Opzione 1: imparare a ricostruire (autoencoder)}
\begin{itemize}
  \item L'\textbf{encoder} estrae le feature rilevanti: codifica l'oggetto nello spazio degli embedding.
  \item Il \textbf{decoder} usa gli embedding per risolvere un compito, es.\ ricostruire (o togliere la sfocatura da) l'immagine originale.
  \item Durante il training è un problema di \textbf{regressione} (predire le componenti RGB dei pixel): \textbf{nessuna label}.
  \item Al test: \textbf{encoder congelato} + una piccola head addestrata con poche label. Ci si aspetta che una specie mai vista (uno scoiattolo
        di un altro colore) finisca vicina agli scoiattoli nello spazio degli embedding, così la head lo classifica bene.
\end{itemize}
L'autoencoder è ripreso in dettaglio per l'anomaly detection nel cap.~\ref{cap:anomaly}.

\subsection{Opzione 2: confronto tra viste (contrastive learning)}
\begin{itemize}
  \item Il soggetto di una foto ruotata, sfocata o decolorata resta lo stesso: è la \textbf{data augmentation}.
  \item Dall'immagine originale (\textbf{anchor}) si creano varianti (\textbf{positivi}): i loro embedding \textbf{devono essere vicini}.
        Un'immagine diversa (\textbf{negativo}) deve avere un embedding \textbf{lontano}.
  \item La loss che premia vicinanza dei positivi e lontananza dei negativi è la \textbf{contrastive loss}: un'altra loss da aggiungere a MSE,
        BCE e cross-entropy. Si chiama ``contrastiva'' perché contrappone il negativo all'anchor e alle sue perturbazioni.
  \item Il feature extractor, poiché codifica l'oggetto nello spazio degli embedding, si chiama anche \textbf{encoder} (slide 4, p.~50).
  \item Utile anche quando i campioni sono pochi (es.\ anomaly detection).
\end{itemize}
\begin{attenzione}[le augmentation introducono bias]
Scegliere le trasformazioni significa iniettare conoscenza: se dico al modello che una foto decolorata è ``la stessa'' immagine, gli insegno
che il colore non conta. Allora riconoscerebbe come falsa la foto di uno \textbf{scoiattolo verde}? Probabilmente no. (Domanda aperta della
slide 4, p.~51; la risposta è un ragionamento, non una regola.)
\end{attenzione}

\begin{sintesi}
\begin{itemize}
  \item Feature categoriche \textbf{ordinali}: ordina con la conoscenza del dominio e mappa su ID sequenziali.
  \item Feature \textbf{nominali}: ID sequenziale (introduce un ordine falso, bias), feature engineering (chiave di lettura che può essere
        sbagliata o invecchiare), one-hot (nessun bias ma alta dimensionalità e vocabolario rigido), \textbf{embedding} appresi (densi,
        $H \ll V$, semantici; solo reti neurali, serve un task).
  \item Vocabolario costruito \textbf{solo sul training}; \code{<UNK>} per i valori mai visti; \code{padding\_idx} per il padding.
  \item \code{nn.Embedding(V, H)} = one-hot + layer lineare: tabella addestrabile in backpropagation, combinabile con FFNN, RNN, GNN.
  \item Embedding \textbf{supervisionati}: ottimi per il task, sub-ottimi per altri, a rischio di scorciatoie (husky vs lupo).
  \item Embedding \textbf{self-supervised} (pretext + downstream): ricostruzione (autoencoder) o contrastive learning (anchor, positivi,
        negativi); le augmentation scelte iniettano bias.
\end{itemize}
\end{sintesi}
